\documentclass[%
 reprint,
superscriptaddress,
nofootinbib,
 amsmath,amssymb,
 aps,
prx,
floatfix,
]{revtex4-2}
\usepackage[hidelinks]{hyperref}
\usepackage{amsmath,amsfonts,amssymb}
\usepackage{xcolor}
\usepackage{graphicx}
\usepackage{float}
\usepackage{amsthm}
\usepackage{appendix}
\usepackage{blkarray}
\usepackage{multirow}
\usepackage{changes}
\usepackage[ruled,vlined]{algorithm2e}

\theoremstyle{definition}
\newtheorem{definition}{Definition}[section]
\newtheorem{theorem}{Theorem}[section]
\newtheorem{lemma}{Lemma}[section]
\newtheorem{proposition}{Proposition}[section]
\newtheorem{remark}[section]{Remark}
\newtheorem{corollary}{Corollary}[section]

\newtheorem{conjecture}[theorem]{Conjecture}
\newtheorem{example}{Example}[section]
\numberwithin{equation}{section}
\newenvironment{hproof}{%
  \renewcommand{\proofname}{ Proof}\proof}{\endproof}
  
\newcommand{\CC}[1]{\textcolor{blue}{(CC) #1}}
\newcommand{\bcl}[1]{\textcolor{orange}{(BCL) #1}}
\DeclareMathOperator{\Tr}{Tr}

\begin{document}

\title{Stabilizer Codes Have Trivial Area Operators}


\author{ChunJun Cao}
 \email{cjcao@vt.edu}
  \affiliation{Joint Center for Quantum Information and Computer Science, University of Maryland, College Park, MD, 20742, USA}
 \affiliation{Institute for Quantum Information and Matter
  California Institute of Technology,
  1200 E California Blvd, Pasadena, CA 91125, USA.}
 \affiliation{Department of Physics, Virginia Tech, Blacksburg, VA, 24061, USA}


\begin{abstract}
We show that no stabilizer codes over any local dimension can support a non-trivial area operator for any bipartition of the physical degrees of freedom even if certain code subalgebras contain non-trivial centers. This conclusion also extends to more general quantum codes whose logical operators satisfy certain factorization properties. The results suggest that non-local ``magic'' would play an important role in reproducing features of gravitational back-reaction and the quantum extremal surface formula. 
We comment on conditions needed to circumvent the no-go result and provide a simple prescription for building non-stabilizer codes that do have non-trivial area operators. We then construct a simple toy example.
\end{abstract}


\maketitle
%\tableofcontents

\section{Introduction}
The intriguing connection between AdS/CFT and quantum error correcting code (QECC) has led to numerous developments in quantum gravity and quantum information\cite{Almheiri:2014lwa}. A key property in AdS/CFT reproduced in QECC models is the Ryu-Takayanagi (RT) formula \cite{Ryu:2006bv,FLM,Hubeny:2007xt}.  In a seminal work \cite{Harlow:2016vwg}, Harlow shows that all quantum erasure correction codes with encoding map $V:\mathcal{H}_L\rightarrow \mathcal{H}$ satisfying complementary recovery produce an RT-like relation such that for any logical state $\rho\in L(\mathcal{H}_L)$, a finite dimensional von Neumann algebra $M_L$ on $\mathcal{H}_L$ and a factorization $\mathcal{H}=\mathcal{H}_A\otimes \mathcal{H}_{\bar{A}}$,
\begin{equation}
    S(\Tr_{\bar{A}} [V\rho V^{\dagger}]) = S(M_L; \rho) +\Tr[\mathcal{L}_A\rho],
    \label{eqn:RT}
\end{equation}
where $\mathcal{L}_A\in L(\mathcal{H}_L)$ is defined to be the \emph{area operator}. 
Analogous to \cite{FLM},  $S(\Tr_{\bar{A}} [V\rho V^{\dagger}])$ corresponds to the entropy of a CFT subregion $A$, $S(M_L; \rho)$ plays the role of the bulk entropy associated with the entanglement wedge of $A$ and $\Tr[\mathcal{L}_A\rho]$ corresponds to the area of the extremal surface  homologous to $A$. In particular, it is argued \cite{Engelhardt:2014gca,Almheiri:2014lwa,Harlow:2016vwg,Cao:2017hrv,Akers:2018fow,Dong:2018seb,Caonote} that in order to emerge effects analogous to gravity and to describe smooth geometries, $\mathcal{L}_A$ needs to be non-trivial, i.e., not proportional to the logical identity.
Although not all quantum codes have a natural geometric description, where the notion of ``area'' is more form than substance, it is crucial to understand what in quantum codes gives rise to non-trivial $\mathcal{L}_A$, and by extension, provides a necessary condition in emerging spacetime and gravity from quantum information. We address this question in our work. 


Popular bottom-up constructions of holographic codes\cite{Pastawski:2015qua,Harris:2018jfl,Jahn:2019nmz,Steinberg:2023wll} have concentrated on the stabilizer formalism. It is known \cite{Harlow:2016vwg,Donnelly:2016qqt,Cao:2020ksw,Pollack:2021yij} that stabilizer codes can admit non-trivial centers in their code subalgebra, of which $\mathcal{L}_A$ is an element. Therefore, such codes may have the capacity to accommodate non-trivial area operators. However, all of the examples that successfully constructed non-trivial area operators are, in fact, non-stabilizer codes \cite{Cao:2020ksw,Pollack:2021yij}. This raises the question whether there are fundamental obstructions in stabilizer codes to having non-trivial $\mathcal{L}_A$. Here we formalize the conjecture by \cite{Pollack:2021yij} and show that no stabilizer code can have non-trivial area operators even if their code subalgebras have non-trivial centers. This no-go result also generalizes to some classes of non-stabilizer codes. These findings further necessitate the importance of long range non-stabilizerness\cite{NLmagic,White:2020zoz,Bao:2022mkc,Liu:2020yso} in understanding the emergence of gravitational features. 


We refer the readers to \cite{Harlow:2016vwg} and \cite{Pollack:2021yij} for details on the RT formula in QECCs. In this article, we only work with codes and von Neumann algebras on finite dimensional Hilbert spaces.  Note that the term ``stabilizer code'' commonly refers to a Pauli stabilizer code in literature where the code subspace is defined by the simultaneous $+1$ eigenspace of all elements in an abelian subgroup of the Pauli group over $n$ qudits with local dimension $q$. Here we adopt the same convention; as such, non-stabilizer codes in this work also include codes like \cite{XS,XP} which extend the Pauli stabilizer formalism\footnote{Strictly speaking, they are also stabilizer codes except the stabilizer group is no longer an abelian Pauli subgroup. However, their codewords in the computational basis need not be stabilizer states.}.





\section{No-go Theorems}
\begin{theorem}\label{thm:1}
Let $\mathcal{C}\subseteq \mathcal{H}$ be a stabilizer code with $M_{\mathcal{C}}$ the associated von Neumann algebra on $\mathcal{C}$. Any area operator ${\mathcal{L}}_{A}\in M_{\mathcal{C}}$ defined in this code must be proportional to the identity.
\end{theorem}
\begin{proof}
Consider any bipartition  of the physical degrees of freedom such that $\mathcal{H}=\mathcal{H}_A\otimes \mathcal{H}_{\bar{A}}$ and let $M_A$ be the code subalgebra supported on $A$. Let $\{|\widetilde{\alpha,ij}\rangle\}$ be a set of orthonormal basis  induced by $M_A$ in the Wedderburn decomposition, which spans $\mathcal{C}=\bigoplus_{\alpha}\mathcal{H}_{a^{\alpha}}\otimes \mathcal{H}_{\bar{a}^{\alpha}}$ such that $\forall\tilde{O}_A\in M_A$, $\tilde{O}_A=\bigoplus_{\alpha}O_{a^{\alpha}}\otimes I_{\bar{a}^{\alpha}}$ , and for operators in the commutant $\forall \tilde{O}'_A\in M'_{{A}}$, $\tilde{O}'_A=\bigoplus_{\alpha}I_{{a}^{\alpha}}\otimes O_{\bar{a}^{\alpha}}$. Here we drop the null space for simplicity.

Recall from \cite{Pollack:2021yij} Lemma C.1 that all stabilizer codes satisfy complementary recovery, that is, one can identify a decomposition $\mathcal{H}_A=\bigoplus_{\alpha}\mathcal{H}_{A_1^{\alpha}}\otimes\mathcal{H}_{A_2^{\alpha}}$,$\mathcal{H}_{\bar{A}}=\bigoplus_{\alpha}\mathcal{H}_{\bar{A}_1^{\alpha}}\otimes\mathcal{H}_{\bar{A}_2^{\alpha}}$
and unitary operators $U_A,U_{\bar{A}}$ such that
\begin{equation}
    |\widetilde{\alpha,ij}\rangle = U_AU_{\bar{A}} \Big(|\alpha,i\rangle_{A_1^{\alpha}}|\alpha,j\rangle_{\bar{A}_1^{\alpha}}|\chi_{\alpha}\rangle_{A_2^{\alpha}\bar{A}_2^{\alpha}}\Big)
\end{equation} 
for some state $|\chi_{\alpha}\rangle$ that captures the residual entanglement between $A$ and $\bar{A}$ which we would then associate with the area of the RT surface.
By the definition of complementary recovery\cite{Harlow:2016vwg}, we know that any element of $M_A'$ is supported on $\bar{A}$, i.e., $M_A'\subseteq M_{\bar{A}}$. Since any operator supported on $\bar{A}$ must commute with those on $A$, we have $M_{\bar{A}}\subseteq M_A'$. It follows that $M_A'=M_{\bar{A}}$.

Let $\chi_{\alpha}=\Tr_{\bar{A}_2^{\alpha}}[|\chi_{\alpha}\rangle\langle\chi_{\alpha}|]$, and the area operator is defined as $\mathcal{L}_A=\bigoplus_{\alpha}S(\chi_{\alpha})\tilde{I}_{a^{\alpha}\bar{a}^{\alpha}}\in Z(M_A)$.
If the center $Z(M_A)$ is trivial, i.e. contains only elements proportional to the logical identity, then so is the area operator. If the center is non-trivial, (e.g.\cite{Donnelly:2016qqt,Cao:2020ksw,Pollack:2021yij}) then 
it can be expressed in the block diagonal form with distinct values of $\alpha$. Each such block is often referred to as an $\alpha$-block in relevant literature in holography.  
We now show that the eigenvalues $S(\chi_{\alpha})$ must be identical for all $\alpha$'s in stabilizer codes. 


Let $\tilde{P}$ be a logical Pauli operator such that $\tilde{P}\in M_{\mathcal{C}}$ and $\Pi_c[\tilde{P},\mathcal{L}_{A}]\Pi_c\ne 0$, where $\Pi_c$ is the projector onto the code subspace $\mathcal{C}$. Note that the logical Pauli operators form a complete operator basis for $M_{\mathcal{C}}$ in a stabilizer $[[n,k,d]]_q$ code over any local dimension $q$. If no such $\tilde{P}$ exists, then $\mathcal{L}_A\in Z(M_{\mathcal{C}})$, leading to a trivial area operator. 


When such a $\tilde{P}$ can be found, let us first decompose the non-trivial area operator as
 \begin{equation}
     \mathcal{L}_A = c_0 \tilde{I}+\sum_i c_i \tilde{Q}_i,
 \end{equation}
 where there is at least one $\tilde{Q}_i$ for which $c_i\ne 0$. 
 Because $\mathcal{L}_A\in Z(M_A)=M_A\cap M_A'$ such $\tilde{Q}_i$ must admit a form that only has support\footnote{This is because each logical Pauli operator is represented as a distinct physical Pauli operator in a stabilizer code. If some of the $\tilde{Q}_i$ have support on $\bar{A}$ also, then their linear combination must be supported on $\bar{A}$. The same need not be true for other more general distinct unitary operators. For instance, let $U_1=(I+iX)/2, U_2=(I-iX)/2$. It is clear that $U_1+U_2=I$ remains unitary, but has trivial support over the qubit.} on $A$.  Thus applying a similar argument to $\mathcal{L}_A\in M_{A'}\subseteq M_{\bar{A}}$, we conclude that $\tilde{Q}_i$ must also admit representations that only have support on $\bar{A}$. This implies that $\tilde{Q}_i$ must be in the intersection of $M_A$ and $M_A'$ hence in $Z(M_A)$. Therefore $[\tilde{Q}_i,\mathcal{L}_A]=0$ and they can be simultaneously diagonalized. Let $\{|\tilde{\psi}_{\alpha}\rangle\}$ be such an eigenbasis over $\mathcal{C}$ for the Pauli operator $\tilde{Q}_i$ where each basis state lives in a particular $\alpha$ block. Note that each such a basis state $\tilde{\rho}^{\alpha}=|\tilde{\psi}_{\alpha}\rangle\langle\tilde{\psi}_{\alpha}|$ can be obtained from unitarily encoding quantum information in the $k$ data qudits (denoted as subsystem $D$) over $n$ physical qudits, such that $\tilde{\rho}^{\alpha}= U_0 (|\psi_{\alpha}\rangle\langle\psi_{\alpha}|_D\otimes |0\rangle\langle 0|_S^{\otimes(n-k)}) U^{\dagger}_0$ where $U_0$ is a Clifford unitary that acts on $\mathcal{H}$. The subsystem $S$ of $n-k$ ancillae are sometimes called the syndrome qudits. For codes that are complementary recoverable, there exist recovery unitary $U_A, U_{\bar{A}}$ such that, 

\begin{align*}
    &U_A U_{\bar{A}}\tilde{P}\tilde{\rho}^{\alpha}\tilde{P}^{\dagger} U_A^{\dagger}U_{\bar{A}}^{\dagger}\\
    = &(U_A U_{\bar{A}} U_0 P_{D}\otimes I_{S}) |\psi_{\alpha}\rangle\langle\psi_{\alpha}|_D\otimes |0\rangle\langle 0|^{\otimes (n-k)}_S  \\
    &\times(U_A U_{\bar{A}} U_0 P_{D}\otimes I_{S})^{\dagger}\\
    = &P' U_A U_{\bar{A}}\tilde{\rho}^{\alpha} ( U_A U_{\bar{A}} P')^{\dagger}\\
    = &P' \rho_{A_1^{\alpha}\bar{A}_1^{\alpha}} \otimes \chi_{A_2^{\alpha} \bar{A}_2^{\alpha}} P'^{\dagger},
\end{align*}

where $P'$ is a Pauli operator. Here we have used the fact that $U_0,U_A,U_{\bar{A}}$ are Clifford unitaries, and hence maps Pauli operators $P_{D}\otimes I_{S}$ to another Pauli $P'$
$$U_AU_{\bar{A}} U_0 P_{D}\otimes I_{S}(U_AU_{\bar{A}} U_0)^{\dagger} = P'.$$
$P_{D}$ acts only on the data qudits in the state $|\psi_{\alpha}\rangle$ while $I_{S}$ acts only on the syndrome qudits (prior to encoding).

 
 Because each state $|\tilde{\psi}_{\alpha}\rangle$ is an eigenstate of a Pauli operator, by Lemma \ref{lemma:1}, we know that $\tilde{P}|\tilde{\psi}_{\alpha}\rangle=|\tilde{\psi}'_{\alpha'}\rangle$ where $|\tilde{\psi}'_{\alpha'}\rangle$ is another eigenstate of $\mathcal{L}_A$. Therefore, we must have 
 
 \begin{equation}
     P' \rho_{A_1^{\alpha}\bar{A}_1^{\alpha}} \otimes \chi_{A_2^{\alpha} \bar{A}_2^{\alpha}} P'^{\dagger} = \rho'_{A_1^{\alpha'}\bar{A}_1^{\alpha'}}\otimes \chi'_{A_2^{\alpha'}\bar{A}_2^{\alpha'}},
     \label{eqn:Pprime}
 \end{equation}
 such that we are taken to a potentially different state on an alpha block with potentially a different eigenvalue. However, note that 
 
 $$P'\rho_{A_1^{\alpha}\bar{A}_1^{\alpha}}\otimes \chi_{A_2^{\alpha}\bar{A}_2^{\alpha}} P'^{\dagger} = P'_1\rho_{A_1^{\alpha}\bar{A}_1^{\alpha}}P'^{\dagger}_1\otimes P'_2 \chi_{A_2^{\alpha}\bar{A}_2^{\alpha}} P'^{\dagger}_2,$$
 
 where both $P'_1, P'_2$ are tensor products of Pauli operators. 
It is obvious that the local unitary $P'_2$ does not alter the entanglement of $\chi_{\alpha}$ across $A_2^{\alpha},\bar{A}_2^{\alpha}$, i.e.,
\begin{align*}
S(\chi'_{\alpha'})&=S(\Tr_{\bar{A}_2}[P'_2\chi_{A_2^{\alpha}\bar{A}_2^{\alpha}}P_2'^{\dagger}]) \\
&= S(P''_2\Tr_{\bar{A}_2}[\chi_{A_2^{\alpha}\bar{A}_2^{\alpha}}]P_2''^{\dagger})=S(\chi_{\alpha}),
\end{align*}
where $P_2''$ is a Pauli substring on $A_2^{\alpha}$ and the entanglement entropy is preserved under unitary conjugation.
As this holds for all of the above basis states, it follows that for any state $\tilde{\rho}\in L(\mathcal{C})$, $\Tr[\tilde{\rho}\mathcal{L}_A]=\sum_{\alpha}p_{\alpha}S(\chi_{\alpha})$ is conserved under action generated by such $\tilde{P}$'s. In other words,
\begin{equation}
    \Pi_c[\mathcal{L}_A,\tilde{P}] \Pi_c=0,
\end{equation}
which is a contradiction. Therefore the area operator must be trivial.
\end{proof}

The proof relies on two crucial elements --- that $Z(M_{\mathcal{C}})$ is trivial, as implied by the definition of the code,  and that Pauli operators are mapped to Pauli operators by Clifford unitaries. 
Note that $Z(M_{\mathcal{C}})$ can be trivial while $Z(M_A)$, the center of the code subalgebra recoverable on $A$, need not be. This is the case, for example, for the $[[4,1,2]]$ code \cite{Cao:2020ksw}. 
Intuitively, 
because recovery unitaries $U_{A},U_{\bar{A}}$ are Cliffords in a stabilizer code, which map logical Pauli operators to Pauli operators, their resulting actions on $\chi_{A_2^{\alpha}\bar{A}_2^{\alpha}}$ do not change the entanglement across  $A_2^{\alpha}$ and $\bar{A}_2^{\alpha}$, which is what gives rise to the RT contribution in the entropy formula (\ref{eqn:RT}). Such conditions are generally broken by non-(Pauli)-stabilizer codes, where $M_{\mathcal{C}}$ can take on more general forms and the decoding unitaries need not map tensor product of operators like $P_{D}\otimes I_{S}$ to tensor product of operators. 

At the same time, we note that such a no-go result continues to hold for a more general class of codes as long as the requisite logical Pauli operator $\tilde{P}$ is mapped to an operator that preserve the entanglement across $A_2^{\alpha}$ and $\bar{A}_2^{\alpha}$, for which the tensor product of Pauli's is sufficient but not necessary. For example, 
\begin{corollary}\label{coro:1}
    Any quantum code that is local unitary (LU) equivalent to a stabilizer code, i.e., $\Pi = U \Pi_c U^{\dagger}$ where $U=\bigotimes_{i=1}^n U_i$, has a trivial area operator.
\end{corollary}

\begin{proof}
    This is obvious because local unitaries applied to each code word does not alter the entanglement structure across any bipartition. Hence for any RT formula $\Tr[\mathcal{L}_A\tilde{\rho_a}]= S(\tilde{\rho_A})-S(M_A;\tilde{\rho}_a)$ is identical to that of the stabilizer code $\rho_a$, meaning that it cannot depend on the codeword in consideration, i.e., $\mathcal{L}_A$ must be proportional to the logical identity.
\end{proof}

Such codes are trivial extensions of the stabilizer code via local unitary deformations, which have local magic, but not non-local magic, i.e., non-stabilizerness that cannot be removed by acting local unitary gates.

Beyond these trivial extensions, there can also be more general codes that have trivial area operators but need not be stabilizer codes as long as certain product structures are preserved under the recovery process. 



\begin{theorem}\label{thm:2}
   For any quantum code $\mathcal{C}$ with an associated code subalgebra $M_{\mathcal{C}}$ satisfying complementary recovery with respect to bipartition $\mathcal{H}=\mathcal{H}_A\otimes\mathcal{H}_{\bar{A}}$,  if $Z(M_{\mathcal{C}})$ is trivial and the logical Pauli operators $\tilde{P}$ are mapped to 
$$U\tilde{P}U^{\dagger}=O_{A_1\bar{A}_1}\otimes O_{A_2}\otimes O_{\bar{A}_2},$$
where $U=U_AU_{\bar{A}}$ are unitary recovery maps, then it admits only a trivial area operator.
\end{theorem}


\begin{proof}
The proof is the same as that for the stabilizer code, except we no longer require $P'$ in (\ref{eqn:Pprime}) to be a Pauli operator. By assumption, it maps $P_{D}\otimes I_{S}$ to $O_{A_1\bar{A}_1}\otimes O_{A_2}\otimes O_{\bar{A}_2}$. This unitary operator does not entangle $A_2$ with $\bar{A}_2$. Hence the rest of the proof follows.  
\end{proof}

\begin{corollary}
If $\mathcal{C}$ is complementary, $Z(M_{\mathcal{C}})$ is trivial and the logical Pauli operators are transversal, then the code has trivial area operators for all bipartitions.
\end{corollary}
\begin{proof}
    If $\tilde{P}$ is transversal, i.e., $\tilde{P}=\bigotimes_i O_i$, then for any bipartition $\mathcal{H}=\mathcal{H}_A\otimes \mathcal{H}_{\bar{A}}$, $U_{A}U_{\bar{A}} \tilde{P} (U_{A}U_{\bar{A}})^{\dagger} = O_{A}\otimes O_{\bar{A}}$. Note that this satisfies the general form required by Theorem~\ref{thm:2}, hence the area operator is trivial. This includes Corollary~\ref{coro:1} as a special case.
\end{proof}
Examples include non-(Pauli-)stabilizer codes such as some XP and XS codes \cite{XP,XS}; hence it is clear that there are non-stabilizer codes (that are not LU equivalent to stabilizer codes) which also have trivial area operators.  Another instance of non-stabilizer code with trivial area operator may be the random tensor networks (RTN) in the large bond dimension limit\cite{Hayden:2016cfa} as they describe fixed area states\cite{Akers:2018fow,Dong:2018seb}. However, RTN models are approximate QECCs and said result holds in the infinite bond dimension limit. Therefore, strictly speaking they are beyond the scope of codes we discuss in this work. 

Interestingly, to produce non-trivial area operators, it appears that we need to consider criteria that are the opposite of what we usually use to construct fault-tolerant gates. Nevertheless, some level of fault-tolerance seems to be retained  in AdS/CFT \cite{threshold} despite this difficulty to incorporate transversal gates. See also \cite{Woods,Faist,Cree:2021rxi}. 

\begin{remark}\label{rmk:erasure}
Any code with a trivial $Z(M_{\mathcal{C}})$ but a non-trivial area operator for all bipartitions 
 $\{A,\bar{A}: |A|,|\bar{A}|\ne 0\}$ cannot have distance $d>1$.
\end{remark}
\begin{proof}
Since area operator is non-trivial, but $Z(M_{\mathcal{C}})$ is, there must exist some logical operator which takes states in one alpha block to another. Let this operator be $\tilde{O}$ and $[\mathcal{L}_A,\tilde{O}]\ne 0$. Recall that $\mathcal{L}_A\in Z(M_A)$, which has support on either $A$ or $\tilde{A}$ entirely, this means that $\tilde{O}$ must have support on both $A$ and $\bar{A}$. If this is true for all bipartitions, then there will always be a subalgebra generated by $\tilde{O}$ that is not recoverable after the erasure of $\bar{A}$. Hence the erasure correction cannot be exact. 
\end{proof}

Although the assumptions of Remark~\ref{rmk:erasure} are somewhat restrictive, the outcome is not terribly surprising in AdS/CFT because any operator that can arbitrarily alter the ``global charge'' like ADM mass should have information leaking to the boundary via soft-hair. Thus any erasure correction is only expected to be approximate at finite $N$. A similar, but less restrictive, distance upper bound is expected in order to reproduce non-trivial boundary 2-point functions\cite{HMERA}. Here the adversarial distance is not a particularly informative measure of the code property because large erasures may still be correctable up to small errors though the recovery is not exact.

It is a common lore that a non-trivial area operator is tied to soft-hair or edge modes such that the ``physical'' Hilbert space in a system with gauge symmetry is not factorizable unlike the logical subspace of a stabilizer code\cite{Buividovich:2008gq,Donnelly:2011hn,Casini:2013rba,Lin:2017uzr}. Such a system can be identified in a stabilizer code by considering a subspace of $\mathcal{C}$ that satisfies additional gauge constraints such that ``emergent gauge symmetry'' is present in a (low energy) subspace $C'$\cite{Levin:2004mi,Dolev:2021ofc,Qi:2022lbd}. Let us examine whether our no-go theorem puts any constraint on such codes. 


For example, what if we do not consider proper stabilizer codes, but an induced version with a more  restricted form of the code subalgebra $\mathcal{N}_A\subset \mathcal{M}_A$? Such an operator algebra QECC generally need not be a stabilizer code. Ref. \cite{Pollack:2021yij} points out that complementary recovery is so restrictive that any von Neumann algebra $M_A$ satisfying such properties must be unique. This is somewhat undesirable because although one may still identify one-sided area operators under $\mathcal{N}_A$ which may be non-trivial, the resulting code cannot satisfy complementary recovery. 

It is also possible to constrain the entire code subalgebra $\mathcal{N}\subset\mathcal{M}$, which amounts to (generalized) concatenation where one effectively defines another code over the $k$ logical qubits/qudits. The no-go theorem does not constrain codes obtained this way, except when the concatenated codes are stabilizer codes. It is known that (generalized) concatenation of stabilizer codes remain stabilizer codes\cite{QL}, hence such combination would yield trivial area operators. For example, by gauging the bulk and constructing a $\mathbb{Z}_2$ gauge theory over $k$ bulk qubits in any holographic stabilizer toy models would still yield trivial area operators. Therefore, while gauging would be useful in generating non-trivial centers, it is insufficient to guarantee the existence of non-trivial ``area operators''. However, it appears that non-local non-stabilizerness, or non-local magic, are essential ingredients in reproducing the desired properties.  



\section{Codes with non-trivial area operator}
We now give an explicit prescription for building codes with non-trivial area operators (Fig.~\ref{fig:circuit}). Each wire $A_1^{\alpha},\bar{A}_1^{\alpha},A_2^{\alpha}$ or $\bar{A}_2^{\alpha}$ in Fig.~\ref{fig:circuit}a is a tensor factor of a Hilbert space which need not have the same dimensions. However, distinct from the usual circuit diagrams, it is understood that these Hilbert spaces are also direct summed over $\alpha$. To distinguish from the usual single and double wires in circuit diagrams that represent tensor factors and classical operators, we use a fuzzy line to denote the implicit direct sum over different $\alpha$ sectors. Solid line on $U_A,U_{\bar{A}}$ denotes action on factorized Hilbert spaces as usual. Circle denotes controlled system. Controlled-controlled gate here is used to indicate the unitary action $$U = \sum_{\alpha}\sum_{i,j} |\alpha;ij\rangle\langle \alpha;ij|\otimes U_{\alpha},$$ which can depend on the $\alpha$-block on which the state of the wire lives. $U_{\alpha}$ acts on the subspace $\mathcal{H}_{A_2^{\alpha}\bar{A}_2^{\alpha}}$. Operationally, it is more convenient to construct such codes with a circuit over tensor factors like physical qubits; we may do so by introducing additional tensor factors to keep track of the $\alpha$ blocks and embedding the circuit into one over a factorizable Hilbert space (Fig.~\ref{fig:circuit}b). Each wire represents a tensor factor, but they need not be of the same dimensionality.

\begin{figure}
    \centering
    \includegraphics[width=0.85\linewidth]{circuit.png}
    \caption{(a) A unitary circuit for preparing a code with non-trivial area operator in an $A$-$\bar{A}$ bipartition. (b) A circuit that achieves the left via embedding into a factorizable Hilbert space. The controlled gate is given by $\sum_{\alpha}|\alpha\rangle\langle\alpha|\otimes U_{\alpha}$. }
    \label{fig:circuit}
\end{figure}




For example, consider the $[[4,1,1]]$ skewed code\cite{Cao:2020ksw} with trivial $A_1,\bar{A}_1$ defined by the encoding isometry $V=|0000\rangle\langle 0|+(|1001\rangle+|0110\rangle)\langle 1|/\sqrt{2}$. The overall encoding tensor is represented as Fig.~\ref{fig:4qubit}b. The codewords $|\tilde{0}\rangle, |\tilde{1}\rangle$ have entanglement structures that are invariant under rotations by construction. Using the protocol above, we can now construct this code with an explicit encoding circuit (Fig.~\ref{fig:4qubit}a). 
\begin{figure}
    \centering
    \includegraphics[width=0.85\linewidth]{4qubit.png}
    \caption{(a) Encoding circuit of the $[[4,1,1]]$ code where the qubits labelled sequentially from top to bottom. The dashed box represent $U_{\alpha}$ while the states $|\alpha\rangle$ correspond to the degrees of freedom in $A^{\alpha}$. (b) The same code represented as an encoding map/tensor with the red line marking the logical input and the black lines marking physical qubits.}
    \label{fig:4qubit}
\end{figure}
By construction, this code has a non-trivial area operator ($S(\chi_{0})=0, S(\chi_{1})=1$) in any bipartition into 2 pairs of adjacent qubits. This is because $\tilde{Z}$ can be represented as any weight-2 $ZZ$ operator that act on adjacent qubits while $\tilde{X}$ acts globally\footnote{Also see example G of \cite{Pollack:2021yij} for an atomic code with non-trivial $A_1,\bar{A}_1$. }. 

Another key feature lies in the non-stabilizerness of the code. This is already apparent in the use of controlled-controlled gates. Let us also study it from the perspective of the Choi state of the encoding channel. Note that the Choi state of a stabilizer code encoding isometry would always have zero magic. The Choi state of the $[[4,1,1]]$ code encoding map is $|V\rangle = V(|00\rangle+|11\rangle)/\sqrt{2}$, which has no local magic as each of its single qubit reduced density matrix is of the form $p|0\rangle\langle 0|+(1-p)|1\rangle\langle 1|$ for some $p>0$. One can further verify  that any 2-qubit reduced density matrix remain stabilizer states. However, the global state has non-trivial non-stabilizerness as measured by the robustness of magic (RoM)\cite{ROM}. 


Unlike the $[[4,1,2]]$ code, when the logical state is not fixed to $|\tilde{1}\rangle$, any single-qubit subsystem is not maximally mixed. We can further convert each 4-qubit state $|\tilde{\psi}\rangle$ to a linear map $V':\mathbb{C}^{2}\rightarrow \mathbb{C}^8$ using the Choi-Jamiolkowski isomorphism. Although state norm is not preserved under $V'$, orthogonality is. A more complicated tensor network encoding map can then be constructed by gluing such atomic ``legos'' with Bell fusion (plus post-selection). However, such tensor networks generally do not represent encoding isometries due to the property of $V'$. Hence they cannot be unitarily constructed. For example, the encoding map $V_2$ in Fig.~\ref{fig:noniso}a encodes two logical qubits, but is non-isometric. Similar to \cite{Cao:2020ksw,Steinberg:2023wll}, operator reconstruction is state-dependent. Explicitly, let $\tilde{X}_1$ be the unitary operator constructed via operator pushing when the logical qubit on the right is fixed to be $|\tilde{1}\rangle$. Note that $\tilde{X}_1$ is a state-specific construction \cite{AP2021} that has support over all 6 qubits. 
\begin{figure}
    \centering
    \includegraphics[width=0.8\linewidth]{noniso.png}
    \caption{(a) Reconstruction of an operator $\tilde{X}_1$ via operator pushing when the second logical qubit is fixed to $|\tilde{1}\rangle$. (b) The holographic toy model using the $[[4,1,1]]$ atomic lego forms a skewed holographic Bacon-Shor code. Fixing logical legs to $|\tilde{0}\rangle$ alters the entanglement structure and the network connectivity by contracting the tensors with $|0\rangle$s (blue triangles). The grey hexagons are perfect tensors\cite{Pastawski:2015qua}.}
    \label{fig:noniso}
\end{figure}
That is, $\tilde{X}_1 V|{\psi}, {1}\rangle = V X_1|\psi, 1\rangle = |\widetilde{(X\psi),{1}}\rangle$. However, the action of $\tilde{X}_1$ on the logical subspace depends on the state. For example, it is easily verified that $\tilde{X}_1V|\psi, 0\rangle \ne VX_1|\psi, 0\rangle = |(\widetilde{X\psi), 0}\rangle$. Hence the reconstruction of this operator is state-dependent. More detailed identification of the code distance, enumerators, operator reconstruction, RoM of Choi states, and the state-dependent reconstruction used above are given in Appendix~\ref{app:1}. 


The symmetries in the tensor make the code a suitable toy example for holographic code constructions. By gluing such atomic codes to a hyperbolic tiling\cite{QL}, the RT surface and area can now can take on a more geometric interpretation given by minimal cuts in the graph. For example, one can construct a more sophisticated model which is a non-stabilizer code but has non-trivial area operator. Fig.~\ref{fig:noniso}b is one such construction, which is an instance of the (gauge-fixed) skewed holographic Bacon-Shor code. For certain bipartitions of the boundary degrees of freedom, the code inherits non-trivial area operators from the $[[4,1,1]]$ legos and the entanglement structure as well as the shape of the RT surface can depend on the bulk state. For more specific properties see \cite{Cao:2020ksw}. 

Notice that each 4 qubit code is a non-isometric map but is orthogonality preserving. This means that the encoding map in Fig.~\ref{fig:noniso}b doesn't have a non-trivial kernel like the non-isometric codes defined in \cite{AP2021,Akers:2022qdl}, but still permits state-dependent reconstructions and causes bulk qubits to ``overlap''\cite{OQ,OQdS,Caonote} in that the commutator of spacelike separated observables is non-vanishing. For this specific code, each bulk state that is fixed to $|\tilde{0}\rangle$ reduces the corresponding tensor to a product state, hence creating a puncture in the network in the order of AdS radius. 

\section{Discussions}
Although stabilizer codes have wide applications in quantum error correction and have reproduced key features of holography, their limitations as toy models are also becoming increasingly apparent. It is known that stabilizer codes cannot reproduce certain features of holography, including the entanglement spectrum (non-flatness), multipartite entanglement\cite{tripartite,markovgap}, correlation functions \cite{Pastawski:2015qua,Jahn:2019nmz,Jahn2022tensornetworkmodels}, and the general expectation of CFTs to have non-local magic\footnote{Note that one can deviate from a stabilizer state on the boundary by encoding a non-stabilizer bulk state. However, the amount of non-stabilizerness injected this way is of $O(1)$ and is insufficient to capture certain holographic features, e.g. the Markov gap \cite{markovgap}, which require non-stabilizerness of $O(1/G_N)$. }. Here we prove that codes without nonlocal magic also cannot capture the features of back-reaction or the quantum extremal surface formula \cite{Engelhardt:2014gca}. 
It further motivates the need to understand holographic magic\cite{White:2020zoz,Goto}, its connection with entanglement, and to construct a more systematic framework for studying non-stabilizer codes.

Given these developments, let us discuss three potential directions to better understand the emergence of gravity in quantum codes, such as the emergence of Einstein's equations (or the Hamiltonian constraint in the absence of dynamics). 
Here we roughly categorize them as 1) the extensions of the stabilizer formalism, 2) gauging, and 3) approximate quantum error correction. 

Although stabilizer codes have many limitations, the general stabilizer formalism remains a powerful machinery \cite{XS,XP,CWS}. Extensions in this front, including considering stabilizer groups with generators that are not given by local tensor products of operators, can be useful for both QECC and holography. Other formalisms, such as operator algebra QEC \cite{OAQEC1,OAQEC2,STABOAQEC}, are also crucial in building codes that circumvent the no-go theorem. Understanding non-trivial centers and area operators in codes with gauge constraints are also important and under-explored. This is the case for ordinary bottom-up QECC models of holography~\cite{Dolev:2021ofc}, which may be advantageous for near-term quantum simulations, but is also true of RTN approaches~\cite{Qi:2022lbd}, which are powerful analytical tools. 
Finally, it is well-accepted that AdS/CFT is best described by approximate QECCs. Therefore, definitions related to the area operator and the RT formula should be modified accordingly.  Ref.\cite{AP2021} made progress in this direction, though instances with non-trivial centers remain to be explored.

\section*{Acknowledgement}
The author thanks Alexander Jahn, Patrick Rall, Jason Pollack, and Yixu Wang for helpful comments and discussions. The author also thanks Ning Bao for (inadvertently) convincing him that it is better to let the note sit on arXiv than in his drawer collecting dust. C.C. acknowledges the support by the Air Force Office of Scientific Research (FA9550-19-1-0360), and the National Science Foundation (PHY-1733907). The Institute for Quantum Information and Matter is an NSF Physics Frontiers Center.




\bibliographystyle{unsrt}
\bibliography{ref} 

\appendix
\section{Further Details}\label{app:1}
We first prove the lemma used in Theorem~\ref{thm:1}.
\begin{lemma}\label{lemma:1}
(Generalized) Pauli operators only changes one eigenstate $|\alpha\rangle$ of a Pauli operator to another $|\alpha'\rangle$ up to a phase. 
\end{lemma}

\begin{proof}
Let $P,Q$ be (generalized) Pauli operators generated by the standard $q\times q$ shift $X$ and clock $Z$ matrices such that they act on the computational basis states $|j\rangle$ as

\begin{align*}
    X|j\rangle &= |j+1~\mathrm{mod}~q\rangle\\
    Z|j\rangle &= \omega^j |j\rangle,
\end{align*}
where $\omega = \exp(2\pi i/q)$.

Suppose $|\phi\rangle$ be an eigenstate of $P$, such that $P|\phi\rangle = \gamma |\phi\rangle$. Now for any Pauli operator $Q$

\begin{equation}
    Q|\phi\rangle = QP^{\dagger} P|\phi\rangle =\gamma QP^{\dagger}|\phi\rangle  
\end{equation}
If $[Q,P^{\dagger}]=0$, then they share the same eigenbasis, i.e., $Q|\phi\rangle\propto|\phi\rangle$.

If $QP=e^{i\omega}PQ$, then $Q|\phi\rangle=\gamma QP^{\dagger}|\phi\rangle = \gamma \exp(-i\omega) P^{\dagger}Q|\phi\rangle$. Hence $Q|\phi\rangle$ is an eigenstate of $P$ with eigenvalue $\gamma^{-1}\exp(i\omega)$. Furthermore, it is clear that $\langle\phi|Q|\phi\rangle = |\gamma|^2 \exp(-i\omega)\langle\phi|Q|\phi\rangle =0$ because the coefficient that is not $1$ has $\omega\ne 0$. 
\end{proof}

The $[[4,1,1]]$ code is complementary in the 2 qubit by 2 qubit bipartition as long as each partition contains adjacent qubits. One can see this by realizing that each subregion reconstructs the $\tilde{Z}$ subalgebra. It is, however, not complementary recoverable in the 3-1 partition as the single qubit does not support a non-trivial logical operation while the remaining 3 qubits do not reconstruct the ${X}$-code subalgebra. In other words, $M_A:=V^{\dagger} L(\mathcal{H}_A)\otimes I_{\bar{A}} V$ is not a von Neumann algebra. However, the area operator can still be defined on at least one side. Note that the codewords are still stabilized by a Pauli subgroup $\langle ZIZI,IZIZ\rangle$, but other symmetries, which are not Pauli operators, exist.

A representation of the logical X operator is
\begin{align*}
    \tilde{X}=\frac{1}{\sqrt{8}}(&IXIX-ZXIX+IXXI+ZXXI\\
    &+XIIX+XZIX-XIXI+XZXI)
\end{align*}
and one can easily read off some representations of $\tilde{Z}$ from the code words as weight-2 $ZZ$ operators acting on adjacent qubits. 

A simple way to verify the distance of the code is through the quantum weight enumerators\cite{SL}. The enumerators of this code are 
\begin{align}
    A(z)=&1+z+2z^2+z^3+3z^4,\\
    B(z)=&1+2z+10z^2+10z^3+9z^4.
\end{align}
Note that $B(z)-A(z)$ has non-zero coefficient at $z^1$, hence $d=1$.
The code is also simple enough that we can verify the distance by observation: the Knill-Laflamme condition is violated for an error of weight 1, i.e., $\langle +|Z_1|-\rangle =1/2\not\propto\delta_{+-}=0.$





The Choi state of this encoding isometry is
$$|V\rangle = \frac{1}{\sqrt{2}}|00000\rangle + \frac 1 2 |10011\rangle + \frac 1 2 |01101\rangle.$$
We compute its robustness of magic $\mathcal{M}(\rho)$ \cite{ROM} for various subsystems. Any 1 or 2-site reduced density matrix is diagonal in the computational basis, and hence are stabilizer states ($\mathcal{M}(\rho_{1,2})=1$). Some subsystems of 3 qubits have non-stabilizerness, though the majority of such subsystems are stabilizer states. The global state has $\mathcal{M}(|V\rangle)=1.7476>1$, hence the state is magical. 


Finally, a state-dependent representation of $\tilde{X}_1\otimes\tilde{I}_2$ for the $[[6,2]]$ code (Fig.~\ref{fig:noniso}a) is
\begin{align*}
  \tilde{X}_1\otimes\tilde{I}_2=& \frac{1}{\sqrt{8}}(IXIXXX-ZXIXXX+IXXIII\\
  &+ZXXIII+XIIXXX+XZIXXX\\
  &-XIXIII+XZXIII).  
\end{align*}
It is easy to verify that it is a unitary operator such that $\tilde{X}_1\otimes\tilde{I}_2|\tilde{01}\rangle=|\tilde{11}\rangle$ but $\tilde{X}_1\otimes\tilde{I}_2|\tilde{00}\rangle\ne|\tilde{10}\rangle$. The overlap with the target state is $\langle \tilde{00}|\tilde{X}_1\otimes\tilde{I}_2|\tilde{10}\rangle=1/\sqrt{2}$.


\end{document}